% VGUFT — Vacuum Grand Unified Fluid Theory
% Monograph scaffold (v0.2 draft). Co-authored in the COMPRESSIBLE site framing
% (vguft.me is canonical; the April incompressible monograph is superseded on
% this point). Every nontrivial claim carries an explicit status tag. The
% established physics being borrowed is tagged [SOLID]; the VGUFT identifications
% are tagged [POSTULATE]; the hard unsolved targets are tagged [OPEN]. Nothing
% here is taken from the unverified third-party "particle-charges" PDF.

\documentclass[11pt]{article}

\usepackage[margin=1in]{geometry}
\usepackage{amsmath,amssymb}
\usepackage[dvipsnames]{xcolor}
\usepackage{hyperref}
\hypersetup{colorlinks=true, linkcolor=RoyalBlue, urlcolor=RoyalBlue, citecolor=RoyalBlue}

% --- VGUFT epistemic status tags ---------------------------------------------
\newcommand{\SOLID}{\textcolor{ForestGreen}{\textbf{[SOLID]}}}
\newcommand{\POSTULATE}{\textcolor{RoyalBlue}{\textbf{[POSTULATE]}}}
\newcommand{\OPENP}{\textcolor{BurntOrange}{\textbf{[OPEN]}}}
\newcommand{\SPEC}{\textcolor{Gray}{\textbf{[SPECULATIVE]}}}

\title{\textbf{VGUFT --- Vacuum Grand Unified Fluid Theory}\\[4pt]
       \large A Dynamical Theory of the Superfluid Vacuum}
\author{Dr.\ Brett Michael Alexander \ (Archmage-Dragon)\\
        \small The Prompt Wizards --- VGUFT Collaboration, Portland, OR}
\date{Version 0.2 (working draft) --- \today}

\begin{document}
\maketitle

\begin{abstract}
VGUFT investigates whether the physical vacuum can be modeled as a
\textbf{compressible, anisotropic, spin-reactive quantum condensate}, and
whether the known forces, particles, and constants emerge from its dynamics.
The medium's transverse (vortical) sector is identified with electromagnetism
and gauge structure; its longitudinal (compressional) sector --- which exists
only because the medium has \emph{finite} compressibility --- is identified with
gravity; and a low-energy acoustic metric supplies relativistic kinematics. The
stance is \emph{emergent, not contradictory}: the substrate must \emph{recover}
general relativity and the Standard Model, not abolish them. This draft states
the governing equations in the compressible framing, borrowing only established
results (each tagged \SOLID), marks the VGUFT identifications as working
hypotheses (\POSTULATE), and is explicit about the unsolved targets (\OPENP).
Adversarial review is welcome; a clean falsification is a successful outcome.
\end{abstract}

\section{Introduction}
The defining question is whether a single dynamical medium can underlie both the
gauge fields and the metric. VGUFT posits one substrate --- a condensate with
density, flow, and internal spin texture --- and treats electromagnetism,
gravity, and relativistic kinematics as distinct \emph{limits} of its dynamics
rather than as independent postulates. The decisive modelling choice, and the
one that distinguishes the current framing from the earlier incompressible
draft, is that the medium is \textbf{compressible}: it is the finite
compressibility that opens a longitudinal channel and therefore a gravitational
sector at all. The framework borrows established physics for every component it
uses (Section~\ref{sec:lineage}) and is explicit about what is assumed versus
what is derived.

\section{The substrate and its order parameter}
The state of the medium is carried by an order parameter with amplitude, phase,
and an internal spin direction,
\begin{equation}
  \Phi^{A}(\mathbf{x},t) \;=\; \sqrt{\rho(\mathbf{x},t)}\;
    e^{\,i\theta(\mathbf{x},t)}\, z^{A}(\mathbf{x},t),
  \qquad z^{\dagger} z = 1 ,
\end{equation}
where $\rho$ is the condensate density (amplitude), $\theta$ the phase (flow
potential), and $z^{A}$ a normalized internal spinor encoding the spin texture.
The superfluid (irrotational) velocity is
\begin{equation}
  \mathbf{v}_s \;=\; \frac{\hbar}{m}\,\nabla\theta .
\end{equation}
The identification of this condensate with the physical vacuum is the central
working hypothesis. \POSTULATE

\section{Governing dynamics (compressible condensate)}
\label{sec:dynamics}
We take the order parameter to obey a Gross--Pitaevskii equation with a
repulsive self-interaction $g>0$ and possible external/self-consistent potential
$V$,
\begin{equation}
  i\hbar\,\partial_t \psi \;=\;
    \Big[-\frac{\hbar^{2}}{2m}\nabla^{2} + g\,|\psi|^{2} + V\Big]\psi ,
  \qquad \psi \equiv \sqrt{\rho}\,e^{i\theta}.
  \label{eq:GP}
\end{equation}
Equation~\eqref{eq:GP} is standard condensate physics. \SOLID\ The Madelung
transform turns it into compressible quantum hydrodynamics: a continuity
equation
\begin{equation}
  \partial_t \rho + \nabla\!\cdot(\rho\,\mathbf{v}_s) = 0 ,
  \label{eq:continuity}
\end{equation}
and a quantum Euler (Bernoulli) equation
\begin{equation}
  m\big(\partial_t \mathbf{v}_s + (\mathbf{v}_s\!\cdot\!\nabla)\mathbf{v}_s\big)
  \;=\; -\,\nabla\!\big(g\rho + V + Q\big),
  \qquad
  Q \equiv -\frac{\hbar^{2}}{2m}\frac{\nabla^{2}\sqrt{\rho}}{\sqrt{\rho}} ,
  \label{eq:euler}
\end{equation}
with $Q$ the quantum potential. \SOLID\ Linearizing
\eqref{eq:continuity}--\eqref{eq:euler} about a uniform background $\rho_0$ gives
Bogoliubov sound with speed
\begin{equation}
  c_s^{2} \;=\; \frac{g\,\rho_0}{m}
  \qquad (\text{long-wavelength limit}).
  \label{eq:sound}
\end{equation}
This \emph{finite} $c_s$ is the linchpin of the framing: a strictly
incompressible medium ($c_s\!\to\!\infty$) has no longitudinal mode and hence no
gravitational channel (Section~\ref{sec:gravity}). \SOLID\ (established
hydrodynamics); the reading of \eqref{eq:sound} as the origin of an emergent
signal speed is \POSTULATE.

\paragraph{Spin texture and an emergent connection.}
Gradients of the internal spinor define a Mermin--Ho gauge connection
$a_\mu = -\,i\,z^{\dagger}\partial_\mu z$, and the superfluid vorticity acquires
the Mermin--Ho term
\begin{equation}
  (\nabla\times\mathbf{v}_s)_i
  \;=\; \frac{\hbar}{2m}\,\epsilon_{ijk}\,
        \hat{\mathbf{n}}\!\cdot\!(\partial_j\hat{\mathbf{n}}\times\partial_k\hat{\mathbf{n}}),
  \label{eq:merminho}
\end{equation}
i.e.\ vorticity sourced by the skyrmion density of the texture $\hat{\mathbf n}$.
Relation~\eqref{eq:merminho} is established for spinor superfluids
($^3$He--A, spinor BECs). \SOLID\ Its identification with a physical gauge field
is \POSTULATE.

\section{Transverse (vortical) channel: electromagnetism}
Define the vorticity and the Lamb vector,
\begin{equation}
  \boldsymbol{\omega} \equiv \nabla\times\mathbf{v}_s ,
  \qquad
  \mathbf{l} \equiv \boldsymbol{\omega}\times\mathbf{v}_s .
\end{equation}
In the transverse (effectively incompressible) limit the hydrodynamic equations
close into a Maxwell-like system under the correspondence
$\mathbf{B}\leftrightarrow\boldsymbol{\omega}$,
$\mathbf{E}\leftrightarrow-\mathbf{l}$:
\begin{align}
  \nabla\!\cdot\boldsymbol{\omega} &= 0, &
  \partial_t\boldsymbol{\omega} + \nabla\times\mathbf{l} &= \nu\nabla^{2}\boldsymbol{\omega}, \\
  \nabla\!\cdot\mathbf{l} &= -\,\partial_t(\nabla\!\cdot\mathbf{v}_s) + s_\rho, &
  \nabla\times\boldsymbol{\omega} &= \frac{1}{c_s^{2}}\,\partial_t\mathbf{l} + \mathbf{j}_\omega ,
\end{align}
where $\nu$ is a kinematic (dissipative) coefficient and $s_\rho,\mathbf{j}_\omega$
collect compressible/nonlinear source terms. The Navier--Stokes$\,\to\,$Maxwell
correspondence is a rigorous mathematical fact (Marmanis 1998; Kambe 2010).
\SOLID\ The identification of physical electromagnetism with this substrate
sector is \POSTULATE. Circulation is topologically quantized,
\begin{equation}
  \Gamma = \oint_{C}\mathbf{v}_s\!\cdot d\boldsymbol{\ell} = n\,\kappa ,
  \qquad \kappa = \frac{h}{m}, \quad n\in\mathbb{Z},
\end{equation}
which is \SOLID\ for superfluids; reading $n$ as the electric-charge quantum is
\POSTULATE\ (and see Section~\ref{sec:charge}).

\section{Longitudinal (compressional) channel: gravity}
\label{sec:gravity}
Small phase perturbations $\theta_1$ on a background flow $(\rho_0,\mathbf{v}_0)$
obey, directly from \eqref{eq:continuity}--\eqref{eq:euler}, a curved-space wave
equation
\begin{equation}
  \partial_\mu\!\big(\sqrt{-g}\,g^{\mu\nu}\,\partial_\nu\theta_1\big) = 0 ,
\end{equation}
with the \emph{acoustic (effective) metric}
\begin{equation}
  g_{\mu\nu} \;\propto\; \frac{\rho_0}{c_s}
  \begin{pmatrix}
    -\,(c_s^{2}-v_0^{2}) & -\,\mathbf{v}_0^{\top}\\[2pt]
    -\,\mathbf{v}_0 & \mathbb{1}
  \end{pmatrix}.
  \label{eq:acoustic}
\end{equation}
The acoustic-metric construction is established analog gravity (Unruh 1981;
Barcel\'o--Liberati--Visser 2005). \SOLID\ The identification of
\eqref{eq:acoustic} with the physical gravitational metric is \POSTULATE, and
the recovery of general-relativistic \emph{precision} tests (Cassini,
GW170817, perihelion/lensing) from the substrate is \OPENP. The essential point
is structural: \eqref{eq:acoustic} exists only because $c_s$ is finite
(Eq.~\eqref{eq:sound}), so gravity in this framework is a consequence of
compressibility.

\section{Emergent relativity}
Because low-energy excitations propagate on the metric \eqref{eq:acoustic}, they
experience an invariant signal speed, effective horizons, and (near suitable
fixed points) an emergent Lorentz symmetry --- the mechanism demonstrated
constructively in superfluid $^3$He--A (Volovik 2003). \SOLID\ as an analog
result; a full derivation of exact Lorentz invariance for the vacuum substrate,
with the observed precision, is \OPENP. The discipline is absolute: a version
predicting no time dilation would be falsified by GPS before lunch.

\section{Particles as vortex-knot solitons; electric charge}
\label{sec:charge}
Particles are modeled as topological excitations --- knotted vortex solitons of
the coupled $(\rho,\theta,z^{A})$ fields, in the Faddeev--Niemi/Hopfion class.
Electric charge is identified with the \textbf{Hopf (linking) invariant} of the
spin texture,
\begin{equation}
  Q_H \;=\; \frac{1}{(4\pi)^{2}}\int \mathbf{A}\!\cdot\!\mathbf{B}\; d^{3}x ,
  \qquad \mathbf{B}=\nabla\times\mathbf{A},
\end{equation}
with $\mathbf{A}$ the connection of the CP$^1$ texture. Knotted Hopfion solitons
are \SOLID\ (they exist as field configurations); the identification $Q_H\!=\!$
electric charge is a deliberate \POSTULATE\ --- a repair of an earlier
charge-from-helicity proposal that contradicted QED --- and deriving the
observed charge spectrum from it remains \OPENP.

\section{Status and open problems (honest list)}
\begin{itemize}
  \item \textbf{Gravity sector} must reproduce GR precision tests
        (Cassini, GW170817, perihelion/lensing) --- not yet demonstrated. \OPENP
  \item \textbf{Charge mechanism} --- the Hopf/linking identification is
        postulated, not derived. \OPENP\ \POSTULATE
  \item \textbf{Chiral fermions} --- recovering the Standard Model's chiral
        structure faces the Nielsen--Ninomiya obstruction. \OPENP
  \item \textbf{Full gauge group} --- recovering $SU(3)\times SU(2)\times U(1)$
        is not established. \OPENP
  \item \textbf{Emergent Lorentz invariance} --- motivated by analog-gravity
        precedent; a full, precision derivation is a target, not a result. \OPENP
  \item \textbf{Equation of state} --- the map from the microscopic interaction
        $g$ to the observed constants ($c$, $G$, $e$) is unspecified. \OPENP
\end{itemize}
The program will have \emph{earned} the definite article (``\emph{The}
dynamical theory\ldots'') on the day it produces a confirmed, novel,
quantitative prediction. Until then it is \emph{a} theory under construction.

\section*{Lineage and borrowed results}
\label{sec:lineage}
Every mechanism above is imported from established physics: the condensate
hydrodynamics (Madelung; Gross--Pitaevskii), vorticity from texture
(Mermin--Ho), the Navier--Stokes$\to$Maxwell correspondence (Marmanis; Kambe),
the acoustic/effective metric and analog horizons (Unruh; Barcel\'o--Liberati--
Visser; Volovik), and knotted solitons (Faddeev--Niemi). VGUFT's content is the
\emph{claim} that the physical vacuum is one such medium and that these sectors
are the physical forces --- claims held to falsifiable standard, not asserted as
settled.

\begin{thebibliography}{9}
  \bibitem{maxwell1865} J.~C.~Maxwell, \emph{A Dynamical Theory of the
    Electromagnetic Field}, Phil.\ Trans.\ R.\ Soc.\ Lond.\ \textbf{155} (1865).
  \bibitem{madelung} E.~Madelung, \emph{Quantentheorie in hydrodynamischer
    Form}, Z.\ Phys.\ \textbf{40} (1927) 322.
  \bibitem{gross} E.~P.~Gross, \emph{Structure of a quantized vortex in boson
    systems}, Nuovo Cimento \textbf{20} (1961) 454;
    L.~P.~Pitaevskii, Sov.\ Phys.\ JETP \textbf{13} (1961) 451.
  \bibitem{merminho} N.~D.~Mermin and T.-L.~Ho, \emph{Circulation and Angular
    Momentum in the $A$ Phase of Superfluid $^3$He}, Phys.\ Rev.\ Lett.\
    \textbf{36} (1976) 594.
  \bibitem{marmanis} H.~Marmanis, \emph{Analogy between the Navier--Stokes
    equations and Maxwell's equations}, Phys.\ Fluids \textbf{10} (1998) 1428.
  \bibitem{kambe} T.~Kambe, \emph{A new formulation of equations of compressible
    fluids by analogy with Maxwell's equations}, Fluid Dyn.\ Res.\ \textbf{42}
    (2010) 055502.
  \bibitem{unruh} W.~G.~Unruh, \emph{Experimental black-hole evaporation?},
    Phys.\ Rev.\ Lett.\ \textbf{46} (1981) 1351.
  \bibitem{blv} C.~Barcel\'o, S.~Liberati, and M.~Visser, \emph{Analogue
    Gravity}, Living Rev.\ Relativity \textbf{8} (2005) 12.
  \bibitem{volovik} G.~E.~Volovik, \emph{The Universe in a Helium Droplet},
    Oxford University Press (2003).
  \bibitem{faddeevniemi} L.~Faddeev and A.~J.~Niemi, \emph{Stable knot-like
    structures in classical field theory}, Nature \textbf{387} (1997) 58.
\end{thebibliography}

\end{document}
